\section{Nonzero checks on independent-chord covers}\label{sec:weighted}

For a stricter check class the limiting rate is zero. Fix a finite
simple connected graph $H$ on $n$ vertices with a spanning $2$-factor,
that is, disjoint cycles covering all vertices. Choose a spanning
tree and orient its chords. Let $\Gamma$ be the free group on these
chords. Assign reciprocal voltages $g_{uv}=g_{vu}^{-1}$ to adjacent
ordered pairs, with tree voltages $1$ and one independent generator
on each oriented chord. Here $g_{uv}$ is in matrix row $u$, column
$v$; it takes a sheet at type $v$ to the sheet multiplied on the
left by $g_{uv}$ at type $u$.

Fix an embedding and chain as in \eqref{eq:chain} for this abstract
$\Gamma$. Over a characteristic-$p$ field $K$, choose arbitrary
diagonal scalars $d_u$ and nonzero scalars $c_{uv}$ for every ordered
edge. Define
\begin{equation}\label{eq:weighted}
 T_{uu}=d_u,\qquad T_{uv}=c_{uv}g_{uv}\quad(uv\in E(H)),
 \qquad T_{uv}=0\quad\text{otherwise}.
\end{equation}
Reverse coefficients need not agree or be inverses.

\begin{lemma}[A surviving determinant monomial]\label{lem:determinant}
\notformalized
Abelianize the chord generators to independent Laurent variables
$z_1,\ldots,z_s$. The resulting matrix $T_z$ has nonzero determinant
over $K[z_1^{\pm1},\ldots,z_s^{\pm1}]$ for every allowed tuple.
\end{lemma}
\begin{informalproof}
Orient the cycles of a spanning $2$-factor and let $\pi_0$ send each
vertex to its successor. The associated determinant term has nonzero
coefficient $\operatorname{sgn}(\pi_0)\prod_uc_{u,\pi_0(u)}$.
For any determinant permutation, record the net directed use of each
undirected edge. This integer current has divergence zero; the
Laurent exponents are its chord coordinates, with fixed choices of
sign. If a second permutation has the same exponents as $\pi_0$,
the difference of the two currents is supported on the spanning
tree and has divergence zero. Leaf deletion shows that it vanishes.

Every oriented factor edge has net use one. Since a permutation
uses each ordered edge at most once, that current forces use of
the same oriented edge. These edges prescribe one outgoing arrow
at every vertex, so the permutation equals $\pi_0$. The selected
monomial is unique and cannot cancel, in any characteristic.
No factor vertex uses a diagonal, so arbitrary $d_u$ cause no issue.
\end{informalproof}

\begin{theorem}[Uniform scalar rate zero]\label{thm:weighted-zero}
\notformalized
For every $\varepsilon>0$ there is $I(\varepsilon)$ such that for
all $i\ge I(\varepsilon)$, all characteristic-$p$ fields, all nonzero
directed edge weights, and all diagonal values in \eqref{eq:weighted},
\[
 \frac{\dim_K\ker T_{Q_i}}{nq_i}<\varepsilon.
\]
\end{theorem}
\begin{informalproof}
First fix one tuple over $K=\F_{p^a}$. Restrict scalars to $\F_p$
to obtain a fixed matrix $T^{\mathrm{res}}$ of size $an$ over
$\F_p[\Gamma]$. Its abelian specialization over
$\F_p(z_1,\ldots,z_s)$ is invertible. To justify this step, an
$\F_p$-basis of $K$ remains linearly independent over $\F_p(z)$:
clear denominators and compare monomial coefficients. Its
$\F_p(z)$-span is a finite-dimensional subalgebra of the field
$K(z)$, hence is a field, since multiplication by any nonzero element
is injective and therefore surjective. It contains $K$ and
$\F_p(z)$, so it is all of $K(z)$. The specialization is consequently
the restriction of scalars of the invertible map $T_z$ from
\cref{lem:determinant}.

Any factorization of $T^{\mathrm{res}}$ through fewer than $an$
columns would specialize to such a factorization of an invertible
matrix over this rational-function field. Hence its inner rank is
$an$. Equation~\eqref{eq:rank-approx} gives a limiting normalized
nullity of zero. Restriction of scalars preserves that normalized
nullity, so the fixed finite-field tuple has rate tending to zero.

Now form the finite-type parameter algebra
\[
 \mathscr R=\F_p[c_{uv},c_{uv}^{-1}:uv\text{ an ordered edge}]
                  [d_u:u\in V(H)].
\]
Its maps to fields are exactly the permitted tuples. Apply
\cref{thm:uniformization} to the universal matrix with $\beta=0$.
The nonzero constraints survive every residue-field specialization.
\end{informalproof}

\begin{corollary}[Arbitrary widths with a common flag]\label{cor:flag-zero}
\notformalized
Replace the scalars in \eqref{eq:weighted} by internal matrices on
one $w_i$-dimensional space at each type, over any characteristic-$p$
field $K_i$. Assume that all edge and diagonal matrices preserve a
common full flag after algebraic closure, and every edge matrix is
invertible. Then the normalized nullity tends to zero uniformly in
the fields, coefficients, flags, and widths $w_i\ge1$.
\end{corollary}
\begin{informalproof}
On each one-dimensional internal flag quotient, the edge matrices
induce nonzero scalars: an invertible triangular matrix has nonzero
diagonal entries. The induced diagonal scalars are unrestricted.
There are $w_i$ scalar check layers, each on $nq_i$ coordinates.
The same threshold in \cref{thm:weighted-zero} bounds every layer's
nullity by $\varepsilon nq_i$. Equation~\eqref{eq:filtered-kernel}
then bounds the full nullity by $\varepsilon nw_iq_i$.
Extending fields preserves rank, so the estimate holds over $K_i$.
\end{informalproof}

Commuting internal edge and diagonal maps give the flag by the proof
of \cref{cor:flag-half}; edge invertibility is still required.
For a recovery check, choose diagonal $I$ and edge coefficients
equal to the negatives of the recovery maps. Thus this theorem
excludes the corresponding full-neighbor, common-flag construction
class, even at unrestricted vector width. It is a bound on specified
checks, not the storage capacity of the underlying graphs.

The nonzero premise cannot be dropped. Set every diagonal to $1$,
retain only one edge $uv$ with coefficients $-1$ in both directions,
and set all other edge coefficients to zero. Its two-type block is
\[
 \begin{pmatrix}I&-P\\-P^{-1}&I\end{pmatrix},
 \qquad P=\rho_{Q_i}(g_{uv}),
\]
which has kernel dimension $q_i$, while all other types have identity
checks. Its normalized nullity is $1/n$ at every level.
No paired-distance hypothesis was needed for the rate-zero theorem;
the independent chord variables and nonzero weights supplied the
stronger rank certificate directly.
